%========================================
% Section：[6] Saṅghābhithuti
%========================================
\section*{m6 | Saṅghābhithuti}

\begin{chantsubsection}
  \chantline
    {Handa mayaṃ saṅghābhithutiṃ karoma se}
    {Now let us recite in praise of the Sangha.}

  \chantline
    {Yo so supaṭipanno bhagavato sāvakasaṅgho}
    {The Sangha of the Blessed One’s disciples has practiced well,}

  \chantline
    {Ujupaṭipanno bhagavato sāvakasaṅgho}
    {has practiced uprightly,}

  \chantline
    {Ñāyapaṭipanno bhagavato sāvakasaṅgho}
    {has practiced in the right way,}

  \chantline
    {Sāmīcipaṭipanno bhagavato sāvakasaṅgho}
    {has practiced properly,}

  \chantline
    {Yadidaṃ cattāri purisayugāni aṭṭha purisapuggalā}
    {namely, the four pairs, the eight kinds of noble individuals.}

  \chantline
    {Esa bhagavato sāvakasaṅgho}
    {This is the Sangha of the Blessed One’s disciples:}

  \chantline
    {Āhuneyyo pāhuneyyo dakkhiṇeyyo añjalīkaraṇīyo}
    {worthy of gifts, worthy of hospitality, worthy of offerings, worthy of reverential salutation,}

  \chantline
    {Anuttaraṃ puññakkhettaṃ lokassa}
    {the unsurpassed field of merit for the world.}

  \chantline
    {Tamahaṃ saṅghaṃ abhipūjayāmi}
    {To this Sangha I pay deep homage.}

  \chantline
    {Tamahaṃ saṅghaṃ sirasā namāmi \;\;\; --- bow ---}
    {To this Sangha I bow my head.}

\end{chantsubsection}